\clearpage\SourceLeaf{134}%
\pdfbookmark[1]{Über eine Aufgabe der sphärischen Astronomie}{sphaerisch1808}
\fancyhead[L]{\small ÜBER EINE AUFGABE DER SPHÄRISCHEN ASTRONOMIE}
\begin{center}
{\small ÜBER EINE AUFGABE DER SPHÄRISCHEN ASTRONOMIE.\par}
\medskip\hrule\medskip
{\small Monatliche Correspondenz zur Beförderung der Erd- und Himmels-Kunde,\par
herausgegeben vom Freiherrn \textsc{von Zach}. October 1808.\par}
\medskip\hrule
\end{center}\medskip
Die tägliche Bewegung der Himmelskörper bietet eine grosse Mannigfaltigkeit
von Problemen dar, welche die Relationen zwischen Stundenwinkeln,
Höhen, Azimuthen, den Örtern der Himmelskörper und der Polhöhe zum Gegenstande
haben. \textsc{Maupertuis} und andere Astronomen haben sich mit verschiedenen
derselben beschäftigt, die indess von sehr ungleicher und zum Theil von
sehr geringer practischer Brauchbarkeit sind. Man hat dergleichen Aufgaben
besonders Seefahrern und Reisenden oder auch solchen Beobachtern empfohlen,
die nur mit wenigen und minder vollkommnen Werkzeugen versehen sind, um
zu den beiden nothwendigsten Beobachtungen, zur Zeit- und Ortsbestimmung,
zu dienen. Denn für die Astronomen, denen die vortrefflichsten Uhren, Mittags-Fernröhre,
Mauer-Quadranten, Vollkreise und Zenith-Sectoren zu Gebote stehen,
braucht nicht gesorgt zu werden: diese können ihre Zeit zu jeder Stunde
aufs schärfste und bequemste bestimmen; und über die zweckmässigsten Methoden
zur Festsetzung der geographischen Lage ihres Beobachtungsorts ist bei ihnen
keine Frage mehr. Diesen Astronomen bleibt natürlich auch die Bestimmung
der Sternpositionen und Sonnen-Örter allein vorbehalten, und der Beobachter
mit schlechtern Instrumenten wird diese immer von jenen entlehnen, wo
er sie braucht, und nicht aus seinen eignen unvollkommnern Beobachtungen ableiten
wollen. Daher ist also z.\,B. die Aufgabe, aus drei beobachteten Sternhö%
\SourceLeaf{135}%
hen zugleich Polhöhe, Declination und Culminationszeit des Sterns zu bestimmen,
von gar keinem practischen Werthe, den einzigen nicht wohl gedenkbaren
Fall ausgenommen, wo man nur einen noch nicht gut bestimmten Stern zu beobachten
Gelegenheit hätte; jenes Verfahren könnte nur dann erträgliche Resultate
geben, wenn die Beobachtungen sehr weit von einander abständen, und zugleich
sowohl der Gang der Uhr während derselben. als die gemessenen Höhen selbst
sehr genau wären, und selbst dann wird man Polhöhe und Zeit immer \textit{viel} schärfer
aus zweien dieser Beobachtungen ableiten können, wenn man die Declination
des Sterns als gegeben ansieht.

Eine der allernützlichsten Aufgaben für Seefahrer und reisende Beobachter
ist die, aus zwei beobachteten Höhen zweier Sterne, deren Rectascensionen und
Declinationen als gegeben angesehen werden, und dem entsprechenden Zeiten der
Uhr, die entweder nach Sternzeit geht oder deren Gang während der Beobachtungen
als bekannt angenommen wird, den Stand der Uhr und die Polhöhe zu
bestimmen. Hier sind die Sterne und die Höhen ganz willkürlich, und man
hat blos die einzige Bedingung zu beobachten, dass die Verticalkreise, in welchen
die Höhen gemessen sind, im Zenith weder einen zu spitzen, noch zu nahe
an $180^0$ fallenden Winkel machen. Auf diese Weise kann man die beiden Höhen
leicht innerhalb des Zeitraums von einigen Minuten messen, unstreitig ein
höchst wichtiger Umstand sowohl für den Seefahrer, der seinen Platz auf dem
Meere stets schnell verändert, als auch für den reisenden Beobachter zu Lande,
dessen Zeit beschränkt ist, oder der vom Wetter nur auf wenige Minuten begünstigt
wird, oder der sich auf den Gang seiner Uhr nicht lange verlassen kann.
Die sich leicht darbietende directe Auflösung dieser Aufgabe beruht auf der Berechnung
von drei sphärischen Dreiecken und ist freilich etwas weitläufig; man
kann aber die indirecte Methode sehr bequem und geschmeidig machen. Ich behalte
mir vor, auf diesen Gegenstand ein andermal zurückzukommen.

Wenn man sich bei Beobachtung von Sternhöhen eines Reflexions\-werkzeuges
und des künstlichen Horizonts bedient, so findet man es anfangs etwas schwierig,
die beiden Bilder ins Feld zu bringen; man erwirbt sich aber hierin bald
eine Fertigkeit, zumal wenn man sich auf die hellern Sterne einschränkt. Eine
vorläufige nur ganz rohe Berechnung der Höhe (falls man Polhöhe und Stand der
Uhr schon beiläufig kennt) erleichtert die Mühe, und wenn man Gelegenheit hat
von einem Stativ zu beobachten, so wird man nicht nur an Bequemlichkeit sehr
\SourceLeaf{136}%
gewinnen, sondern auch die Berührungen weit schärfer und besser in der Mitte
des Gesichtsfeldes bemerken können, als wenn man aus freier Hand beobachten
muss. Inzwischen durch fleissige Übung wird man es auch hierin weit bringen
können.

Eine Unbequemlichkeit bei der vorhin beschriebenen Methode besteht darin,
dass man des Nachts bei Licht die feinen Abtheilungen des Sextanten nicht so
gut ablesen kann, als bei Tage. Wichtiger ist indessen der erst seit einiger Zeit
zur Sprache gekommene Umstand, dass wenigstens viele auch von den ersten
Meistern verfertigte Spiegel-Sextanten in Ansehung der Theilung nicht ganz den
Grad von Vollkommenheit haben, den man ohne weiteres ihnen zuzutrauen bisher
bei uns gewohnt war. Ich selbst besitze einen 10zolligen von \textsc{Troughton}
Nro.\,420, der zwar übrigens vortrefflich ist, aber ganz entschieden die Winkel
von $100$ bis $120^0$ um $50$ bis $60$ oder $70''$ zu klein gibt. Da übrigens alle Prüfungen
und Berichtigungen auf das sorgfältigste damit vorgenommen sind, so
kann ich dies blos einem Theilungsfehler zuschreiben. Auf dem Meere sind
zwar bei Breiten- und Zeitbestimmungen solche Fehler von gar keiner, und auf
dem festen Lande bei Örtern, deren Lage bis dahin noch ganz unbekannt war,
von geringer Bedeutung; allein in solchen Fällen, wo man grössere Schärfe zu
erreichen wünscht, darf man die Vollkommenheit der Theilung nicht auf Treue
und Glauben annehmen, man muss entweder erst die Theilungsfehler mit möglichster
Genauigkeit auszumitteln suchen, oder darauf Verzicht thun, durch Methoden,
die scharf gemessene Höhen voraussetzen, ganz zuverlässige Bestimmungen
zu machen.

Aus diesen Gründen wird vielen Beobachtern eine Methode nicht unwillkommen
sein, nach der man auch mit einem noch so schlecht oder allenfalls gar
nicht getheilten Instrumente in kurzer Zeit Polhöhe und Stand der Uhr mit
grosser Schärfe bestimmen und so auch die etwaigen Theilungsfehler des Instruments
selbst bestimmen kann. Die Genauigkeit der Resultate hängt hier also
lediglich von der Sorgfalt ab, womit man die Berührung der Bilder beobachtet
hat, indem bei dem heutigen Zustande unserer Sternverzeichnisse die anzuwendenden
Positionen der Sterne als vollkommen fehlerfrei betrachtet werden dürfen.
Diese Methode besteht darin, dass man die Zeiten abwartet, wo drei beliebige
Sterne in Verticalkreisen, die am Zenith nicht zu spitze Winkel machen, einerlei,
übrigens willkürliche, Höhe erreichen, welche selbst nicht bekannt zu sein
\SourceLeaf{137}%
braucht. Die bequemste Art hieraus Polhöhe und Stand der Uhr abzuleiten,
wird folgende sein. Ich bezeichne mit
\begin{description}
\item[$\alpha,\alpha',\alpha''$] die geraden Aufsteigungen der drei Sterne,
\item[$\delta,\delta',\delta''$] die Abweichungen derselben, südlich als negativ betrachtet,
\item[$\theta,\theta',\theta''$] die drei Zeiten an der Uhr, wo diese Sterne die Höhe $h$ erreichen,
\item[$k$] Voreilung der Uhr vor Sternzeit, welche ich für alle drei Beobachtungen
als gleich annehme; liefe die Uhr nicht nach Sternzeit, so könnte
man unter $k$ ihre Voreilung für irgend ein willkürliches Zeitmoment annehmen,
dann würden aber unter $\theta,\theta',\theta''$ die Zeiten zu verstehen sein,
die die Uhr gezeigt haben würde, wenn sie zwischen den Beobachtungen
und jenem Zeitmoment genau Sternzeit gehalten hätte. Man könnte
also etwa unter $k$ die Voreilung bei der ersten Beobachtung verstehen,
und die andern beiden Zeitmomente gehörig vermehren oder vermindern,
wenn die Uhr langsamer oder schneller als Sternzeit liefe.
\item[$\varphi$] die Polhöhe des Ortes.
\end{description}
Offenbar werden also $\theta-k-\alpha$, $\theta'-k-\alpha'$, $\theta''-k-\alpha''$, in Bogen verwandelt,
die drei Stundenwinkel sein, und folglich wird man, wenn man
\[
\begin{aligned}
\theta-\alpha&=t\\
\theta'-\alpha'&=t'\\
\theta''-\alpha''&=t''
\end{aligned}
\]
setzt, folgende drei Gleichungen haben:
\[
\begin{aligned}
\mathrm{I.}\qquad&\sin h=\sin\varphi\sin\delta+\cos\varphi\cos\delta\cos(t-k)\\
\mathrm{II.}\qquad&\sin h=\sin\varphi\sin\delta'+\cos\varphi\cos\delta'\cos(t'-k)\\
\mathrm{III.}\qquad&\sin h=\sin\varphi\sin\delta''+\cos\varphi\cos\delta''\cos(t''-k)
\end{aligned}
\]

Zieht man I von II ab, so wird nach einer leichten Verwandlung
\[
\begin{aligned}
&2\sin\varphi\sin\tfrac12(\delta'-\delta)\cos\tfrac12(\delta'+\delta)\\
&\quad=2\cos\varphi\cos[\tfrac12(t+t')-k]\cos\tfrac12(t'-t)\sin\tfrac12(\delta'-\delta)\sin\tfrac12(\delta'+\delta)\\
&\qquad+2\cos\varphi\sin[\tfrac12(t+t')-k]\sin\tfrac12(t'-t)\cos\tfrac12(\delta'-\delta)\cos\tfrac12(\delta'+\delta)
\end{aligned}
\]
oder
\[
\begin{aligned}
\operatorname{tang}\varphi&=\cos[\tfrac12(t+t')-k]\cos\tfrac12(t'-t)\operatorname{tang}\tfrac12(\delta'+\delta)\\
&\quad+\sin[\tfrac12(t+t')-k]\sin\tfrac12(t'-t)\operatorname{cotg}\tfrac12(\delta'-\delta)
\end{aligned}
\]
\SourceLeaf{138}%
Man bestimme $A'$ und $B'$ so, dass
\[
\begin{aligned}
A'\sin B'&=\sin\tfrac12(t'-t)\operatorname{cotg}\tfrac12(\delta'-\delta)\\
A'\cos B'&=\cos\tfrac12(t'-t)\operatorname{tang}\tfrac12(\delta'+\delta)
\end{aligned}
\]
wird, so verwandelt sich die vorige Gleichung, wenn man zugleich
\[
\tfrac12(t+t')-B'=C'
\]
setzt, in folgende
\[
\mathrm{IV.}\qquad\operatorname{tang}\varphi=A'\cos(C'-k)
\]

Völlig auf gleiche Weise, oder bloss durch Vertauschung der Grössen, die
sich auf die zweite Beobachtung beziehen, mit denen, die sich auf die dritte beziehen,
übersieht man, dass, wenn $A''$ und $B''$ so bestimmt werden, dass
\[
\begin{aligned}
A''\sin B''&=\sin\tfrac12(t''-t)\operatorname{cotg}\tfrac12(\delta''-\delta)\\
A''\cos B''&=\cos\tfrac12(t''-t)\operatorname{tang}\tfrac12(\delta''+\delta)
\end{aligned}
\]
wird, und man zugleich
\[
\tfrac12(t+t'')-B''=C''
\]
setzt, folgende Gleichung sich ergeben wird
\[
\mathrm{V.}\qquad\operatorname{tang}\varphi=A''\cos(C''-k)
\]

Aus der Verbindung von IV und V lässt sich nun leicht $k$ und $\varphi$ ableiten.
Man hat nemlich
\[
A'\cos(C'-k)=A''\cos(C''-k)
\]
folglich
\[
\begin{aligned}
&(A''-A')\cos[\tfrac12(C'+C'')-k]\cos\tfrac12(C''-C')\\
&=(A'+A'')\sin[\tfrac12(C'+C'')-k]\sin\tfrac12(C''-C')
\end{aligned}
\]
Setzt man also
\[
\frac{A'}{A''}=\operatorname{tang}\zeta
\]
wodurch
\[
\frac{A''-A'}{A''+A'}=\operatorname{tang}(45^0-\zeta)
\]
\SourceLeaf{139}%
wird, und bestimmt $\psi$ durch die Gleichung
\[
\operatorname{tang}(45^0-\zeta)\operatorname{cotang}\tfrac12(C''-C')=\operatorname{tang}\psi
\]
so hat man
\[
k=\tfrac12(C'+C'')-\psi
\]
und nachher $\varphi$ durch eine der beiden Gleichungen IV oder V. Aus der gefundenen
Polhöhe und den Stundenwinkeln $t-k$, $t'-k$, $t''-k$ kann man nachher,
wenn man will, durch die Formeln I, II, III oder durch andre bekannte
Methoden die Höhe $h$ ableiten, die aus allen drei Beobachtungen denselben
Werth erhalten muss.

Dieser Auflösung sind noch folgende Bemerkungen beizufügen:

1) Um aus zwei Gleichungen $A\sin B=M$, $A\cos B=N$, die Grössen
$A$ und $B$ zu bestimmen, bedient man sich der Formeln
\[
\begin{gathered}
\operatorname{tang}B=\frac{M}{N}\\
A=\frac{M}{\sin B}\quad\text{oder}\quad A=\frac{N}{\cos B}
\end{gathered}
\]
Den erstern Ausdruck für $A$ zieht man vor, wenn $M$ grösser ist als $N$, im entgegengesetzten
Falle den andern. Bei der Bestimmung des Winkels $B$ durch
seine Tangente bleibt die Wahl zwischen dem ersten und dritten Quadranten, wenn
die Tangente positiv, oder zwischen dem zweiten und vierten, wenn sie negativ
ist, willkürlich; man wählt denjenigen, wo der Sinus das Zeichen von $M$ und
der Cosinus das Zeichen von $N$ hat, wodurch $A$ positiv wird.

2) Die ähnliche Zweideutigkeit bei der Bestimmung von $\psi$ durch die Tangente
muss so entschieden werden, dass $\operatorname{tang}\varphi$ positiv wird; man nimmt also
$\psi$ zwischen $-90^0$ und $+90^0$, vorausgesetzt nemlich, dass die Beobachtungen
in der nördlichen Hemisphäre gemacht sind. In der südlichen wäre es umgekehrt.

Bei allen Methoden, die man dem practischen Astronomen zu seinem Gebrauche
vorschlägt, ist es eine unerlässliche Pflicht, dass man den Einfluss der
unvermeidlichen Beobachtungsfehler auf die Resultate würdige, damit man sich
überzeugen kann, ob sie überhaupt und unter welchen Umständen sie mit Sicherheit
anwendbar sind. Der Vernachlässigung dieser Pflicht hat man die vielen
unreifen Einfälle zuzuschreiben, über deren Unwerth die practischen Astronomen
klagen. Eine leichte, hier aber der Kürze wegen zu unterdrückende Un%
\SourceLeaf{140}%
tersuchung zeigt, dass, wenn $\lambda,\lambda',\lambda''$ die den drei Beobachtungen entsprechenden
Azimuthe sind (vom Südpunkte an nach der Richtung der täglichen Bewegung
gezählt), ein Fehler von $\Delta$ Zeit-Secunden bei der ersten Berührung einen
Fehler von
\[
\frac{\Delta\sin\lambda\sin\tfrac12(\lambda''+\lambda')}{2\sin\tfrac12(\lambda'-\lambda)\sin\tfrac12(\lambda''-\lambda)}\quad\text{Zeit-Secunden}
\]
bei der Zeitbestimmung, und von
\[
\frac{15\Delta\sin\lambda\cos\varphi\cos\tfrac12(\lambda''+\lambda')}{2\sin\tfrac12(\lambda'-\lambda)\sin\tfrac12(\lambda''-\lambda)}\quad\text{Bogen-Secunden}
\]
bei der Polhöhe nach sich zieht; um diese Grössen müssen nemlich $k$ vermindert
und $\varphi$ vermehrt werden, wenn man die Berührung zu früh beobachtet hat.
Die Ausdrücke für den Einfluss der Fehler bei der zweiten und dritten Beobachtung
sind diesen ganz ähnlich und entstehen blos durch Vertauschung von $\lambda$ mit
$\lambda'$ oder $\lambda''$. Da übrigens die Höhe des ersten Sterns in $\Delta$ Zeit-Secunden um
$15\Delta\sin\lambda\cos\varphi$ Raum-Secunden abnimmt, so lässt sich der Einfluss der Beobachtungsfehler
noch einfacher so ausdrücken: Wenn der erste Stern in dem Moment,
wo man seine Höhe mit der, auf welche das Instrument gestellt war, übereinstimmend
fand, wirklich noch oder schon um $D$ Raumsecunden höher war, so
erhält man $k$ in Raum-Secunden zu gross um
\[
\frac{D\sin\tfrac12(\lambda''+\lambda')}{2\cos\varphi\sin\tfrac12(\lambda'-\lambda)\sin\tfrac12(\lambda''-\lambda)}
\]
und $\varphi$ zu klein um
\[
\frac{D\cos\tfrac12(\lambda''+\lambda)}{2\sin\tfrac12(\lambda'-\lambda)\sin\tfrac12(\lambda''-\lambda)}
\]

Hieraus folgt nun, dass man blos dahin zu sehen hat, dass von den Sinussen
von $\tfrac12(\lambda'-\lambda)$, $\tfrac12(\lambda''-\lambda)$, $\tfrac12(\lambda''-\lambda')$ keiner zu klein wird, welches man
dadurch bewirkt, dass man nur Sterne auswählt, die die Höhe $h$ in ziemlich
ungleichen Azimuthen erreichen. Zweitens aber ist klar, dass Sterne, deren
Höhe sich langsam ändert, eben so brauchbar sind, als solche, die schnell steigen
oder fallen; es kommt bei jenen durchaus nicht darauf an, dass man den
Augenblick, wo sie die verlangte Höhe haben, haarscharf trifft, sondern nur,
dass sie in dem Augenblick, den man dafür annimmt, wirklich nicht merklich
davon abstehen. Man kann demnach auch ohne Bedenken Sterne nahe bei der
Culmination oder den Polarstern wählen, und gerade solche sind sehr zweck%
\SourceLeaf{141}%
mässig, weil man da dem eben erwähnten Erforderniss mit Ruhe Genüge thun
kann. Einer von den drei Sternen wenigstens wird übrigens immer seine Höhe
schneller ändern, wenn die Bedingung der ungleichen Azimuthe erfüllt ist.

Ehe man diese Beobachtung vornimmt, ist es gut sich darauf vorzubereiten.
Man wird leicht in jeder Stunde einige kenntliche Sterne auffinden, die
bald nach einander gleiche Höhe erreichen. Blos mit einem Globus oder einer
stereographischen Projection der Himmelskugel wird man dies leicht bewerkstelligen
können. Weiss man die Zeitintervalle, wie die drei Beobachtungen auf
einander folgen, auf ein Paar Minuten voraus, so wird die Beobachtung desto
besser gelingen. Man müsste schlecht beobachtet haben, wenn man so nicht wenigstens
auf ein Paar Secunden den Stand der Uhr erhielte; durch Combination
mehrerer Sterne und durch Wiederholung an mehrern Abenden, wo man auch
dieselben Höhen von neuem messen kann, wird man es in seiner Gewalt haben
vermittelst dieser Methode einen sehr hohen Grad von Genauigkeit zu erreichen.
Übrigens lassen sich besonders für den Fall, wo man dieselben Sterne in derselben
oder fast derselben Höhe öfters beobachtet hat, mehrere Abkürzungen der
Rechnung geben, bei welchen ich mich aber hier nicht aufhalte.

Um diese Methode noch mehr zu erläutern, füge ich die Berechnung für
eine wirkliche Anwendung hier in extenso bei. Den 27. August d. J. beobachtete
ich auf dem Quecksilber-Horizont mit meinem auf die doppelte Höhe
$105^0\ 18'\ 55''$ gestellten Sextanten aus freier Hand die Sterne $\alpha$ Andromeda,
$\alpha$ kleiner Bär, $\alpha$ Leyer an der genau nach Sternzeit gehenden \textsc{Shelton}schen
Pendel-Uhr, wie folgt:
\begin{center}
\begin{tabular}{lrrr}
$\alpha$ Andromeda $\ .\ .\ .\ .$&$21^{\mathrm u}$&$33^{\mathrm m}$&$26^{\mathrm s}$\\
$\alpha$ kleiner Bär $\ .\ .\ .\ .$&$21$&$47$&$30$\\
$\alpha$ Leyer $\ .\ .\ .\ .\ .\ .$&$22$&$5$&$21$
\end{tabular}
\end{center}
Die scheinbaren Stellungen der Sterne an diesen Tagen sind folgende:
\begin{center}
\small
\begin{tabular}{lcc}
&Gerade Aufsteigung&Nördl. Abweichung\\
$\alpha$ Andromeda&$23^{\mathrm u}\ 58^{\mathrm m}\ 33^{\mathrm s}33$&$28^0\ 2'\ 14''8$\\
$\alpha$ kleiner Bär&$0\quad55\quad4.7$&$88\quad17\quad5.7$\\
$\alpha$ Leyer $\ .\ .$&$18\quad30\quad28.96$&$38\quad37\quad6.6$
\end{tabular}
\end{center}
Hieraus wird
\SourceLeaf{142}%
\begin{center}\small
\begin{tabular}{rcrcl}
&in Zeit&&&in Bogen\\
$t=$&$21^{\mathrm u}\ 34^{\mathrm m}\ 52^{\mathrm s}67$&$\ .\ .\ .$&&$323^0\ 43'\ 10''05$\\
$t'=$&$20\quad52\quad25.30$&$\ .\ .\ .$&&$313\quad6\quad19.50$\\
$t''=$&$3\quad34\quad52.04$&$\ .\ .\ .$&&$53\quad43\quad0.60$
\end{tabular}
\par\medskip
\begin{minipage}[t]{.49\textwidth}\centering
\begin{tabular}[t]{@{}lr@{}}
$\tfrac12(t'-t)=$&$-5^0\ 18'\ 25''27$\\
$\tfrac12(t'+t)=$&$318\quad24\quad44.77$\\[6pt]
$\tfrac12(\delta'-\delta)=$&$30\quad7\quad25.45$\\
$\tfrac12(\delta'+\delta)=$&$58\quad9\quad40.25$\\[6pt]
$\log\sin\tfrac12(t'-t)\;.\;.$&$8.9661069\,\mathrm n$\\
$\log\cot\tfrac12(\delta'-\delta)\;.\;.$&$0.2363974$\\[6pt]
$\log\cos\tfrac12(t'-t)\;.\;.$&$9.9981343$\\
$\log\operatorname{tang}\tfrac12(\delta'+\delta)\;.\;.$&$0.2069331$
\end{tabular}
\end{minipage}\hfill
\begin{minipage}[t]{.49\textwidth}\centering
\begin{tabular}[t]{@{}lr@{}}
$\tfrac12(t''-t)=$&$-135^0\ 0'\ 4''72$\\
$\tfrac12(t''+t)=$&$188\quad43\quad5.32$\\[6pt]
$\tfrac12(\delta''-\delta)=$&$5\quad17\quad25.90$\\
$\tfrac12(\delta''+\delta)=$&$33\quad19\quad40.70$\\[6pt]
$\log\sin\tfrac12(t''-t)\;.\;.$&$9.8494751\,\mathrm n$\\
$\log\cot\tfrac12(\delta''-\delta)\;.\;.$&$1.0333869$\\[6pt]
$\log\cos\tfrac12(t''-t)\;.\;.$&$9.8494949\,\mathrm n$\\
$\log\operatorname{tang}\tfrac12(\delta''+\delta)\;.\;.$&$9.8179461$
\end{tabular}
\end{minipage}
\end{center}
Wir haben folglich
\begin{center}\small
\begin{minipage}[t]{.49\textwidth}\centering
\begin{tabular}[t]{@{}lr@{}}
$\log A'\cos B'\;.\;.\;.$&$0.2050674$\\
$\log A'\sin B'\;.\;.\;.$&$9.2025043\,\mathrm n$
\end{tabular}
\end{minipage}\hfill
\begin{minipage}[t]{.49\textwidth}\centering
\begin{tabular}[t]{@{}lr@{}}
$\log A''\cos B''\;.\;.\;.$&$9.6674410\,\mathrm n$\\
$\log A''\sin B'\;.\;.\;.$&$0.8828620\,\mathrm n$
\end{tabular}
\end{minipage}
\end{center}
Woraus wir erhalten
\begin{center}\small
\begin{minipage}[t]{.49\textwidth}\centering
\begin{tabular}[t]{@{}lr@{}}
$B'=$&$354^0\ 19'\ 22''04$\\
$B''=$&$266\quad30\quad55.07$\\[6pt]
$C'=$&$-35\quad54\quad37.27$\\
$C''=$&$-77\quad47\quad49.75$\\[6pt]
$\tfrac12(C''-C')=$&$-20^0\ 56'\ 36''24$\\
$\tfrac12(C''+C')=$&$-56\quad51\quad13.51$
\end{tabular}
\end{minipage}\hfill
\begin{minipage}[t]{.49\textwidth}\centering
\begin{tabular}[t]{@{}lr@{}}
$\log A'\;.\;.\;.$&$0.2072029$\\
$\log A''\;.\;.\;.$&$0.8836657$\\[6pt]
$\log\operatorname{tang}\zeta\;.\;.$&$9.3235372$\\
$\zeta=$&$11^0\ 53'\ 41''28$\\[6pt]
$45^0-\zeta=$&$33\quad6\quad18.72$
\end{tabular}
\end{minipage}
\par\medskip
\begin{tabular}{lr}
$\log\operatorname{tang}(45^0-\zeta)\;.\;.$&$9.8142617$\\
$\log\cot\tfrac12(C''-C')\;.\;.$&$0.4171063\,\mathrm n$\\\hline
$\log\operatorname{tang}\psi\;.\;.\;.\;.$&$0.2313680\,\mathrm n$\\
\multicolumn{2}{l}{$\psi=-59^0\ 35'\ 14''71$}\\
\multicolumn{2}{l}{$k=+2\quad44\quad1.20$}
\end{tabular}
\end{center}
Die Uhr eilte also der Stern-Zeit vor um $10^{\mathrm m}\,56^{\mathrm s}08$, welches auf ein Paar Zehn%
\SourceLeaf{143}%
theile mit dem übereinstimmt, was aus Stern-Durchgängen am Mauer-Quadranten
abgeleitet war.

Zur Bestimmung der Polhöhe haben wir
\[
\begin{aligned}
C'-k&=-38^0\ 38'\ 38''47\\
C''-k&=-80\quad31\quad50.95
\end{aligned}
\]
\begin{center}\small
\begin{minipage}[t]{.49\textwidth}\centering
\begin{tabular}[t]{@{}lr@{}}
$\log A'\;.\;.\;.\;.$&$0.2072029$\\
$\log\cos(C'-k)\;.$&$9.8926738$\\\hline
$\log\operatorname{tang}\varphi\;.\;.\;.$&$0.0998767$\\
\multicolumn{2}{l}{$\varphi=51^0\ 31'\ 51''51$}
\end{tabular}
\end{minipage}\hfill
\begin{minipage}[t]{.49\textwidth}\centering
\begin{tabular}[t]{@{}lr@{}}
$\log A''\;.\;.\;.\;.$&$0.8836657$\\
$\log\cos(C''-k)\;.$&$9.2162109$\\\hline
&$0.0998766$
\end{tabular}
\end{minipage}
\end{center}

Berechnet man mit diesen Resultaten die wahre Höhe der Sterne, so findet
man aus allen drei Sternen auf ein Hunderttheil einer Secunde übereinstimmend
\[
h=52^0\ 37'\ 21''3
\]
Die Refraction, nach Barometer- und Thermometer-Stand verbessert, ist
$42''7$, also scheinbare Höhe $52^0\ 38'\ 4''$. Der Collimationsfehler des Sextanten
ist $-3'\ 30''$, also der gemessene Winkel $105^0\ 15'\ 25''$, welcher also, weil der
wahre $105^0\ 16'\ 8''$ ist, um $43''$ zu klein ist.

Am 25. August war der Sextant auf einen $5''$ grösseren Winkel, nemlich
$105^0\ 19'\ 0''$ gestellt, dieselben drei Sterne wurden bei folgenden Uhrzeiten beobachtet:
\begin{center}
\begin{tabular}{lrrr}
$\alpha$ Andromeda $\ .\ .\ .\ .$&$21^{\mathrm u}$&$33^{\mathrm m}$&$29^{\mathrm s}$\\
$\alpha$ kleiner Bär $\ .\ .\ .\ .$&$21$&$47$&$38$\\
$\alpha$ Leyer $\ .\ .\ .\ .\ .\ .$&$22$&$5$&$22$
\end{tabular}
\end{center}
Hieraus folgt:
\begin{center}
\begin{tabular}{lr}
Voreilung der Uhr vor Sternzeit&$10^{\mathrm m}\,57^{\mathrm s}9$\\
$\varphi=$&$51^0\ 31'\ 56''7$\\
$h=$&$52\quad37\quad29.15$
\end{tabular}
\end{center}
Doppelte scheinbare Höhe $105^0\ 16'\ 25''5$; nach dem Sextanten $105^0\ 15'\ 30''$,
also Fehler des Sextanten $-54''5$.

Im Mittel also: Fehler des Sextanten bei dem Winkel $105^0$, $-48''7$;
Polhöhe $51^0\ 31'\ 54''1$. Einige Beobachtungen des Herrn Prof.\ \textsc{Harding} mit den%
\SourceLeaf{144}%
selben Sternen, aber in einer etwas grössern Höhe, so wie eine frühere Beobachtung von mir mit $\alpha$ Andromeda, $\alpha$ Adler und $\alpha$ Leyer, geben für die Polhöhe bis auf ein Paar Secunden dasselbe Resultat.

Um den Einfluss der Beobachtungsfehler desto besser übersehen zu können, habe ich nach den Zahlen des hier entwickelten Beispiels die den drei Sternen entsprechenden Azimuthe berechnet und gefunden:
\begin{center}
\begin{tabular}{lcl}
$\alpha$ Andromeda & $.\ .\ .\ .$ & $293^{0}\ 45'\ 15''$\\
$\alpha$ kleiner Bär & $.\ .\ .\ .$ & $182\quad9\quad9$\\
$\alpha$ Leyer & $.\ .\ .\ .\ .\ .$ & $90\quad17\quad52$
\end{tabular}
\end{center}
Hieraus finde ich nach obigen Formeln, dass, wenn die drei Berührungen um $\Delta,\Delta',\Delta''$ Zeit-Secunden zu früh beobachtet sind, die Voreilung der Uhr um
\[
+0.391\Delta+0.0066\Delta'+0.603\Delta''
\]
Zeit-Secunden zu klein, und die Polhöhe um
\[
+3.808\Delta-0.2884\Delta'-3.519\Delta''
\]
Raum-Secunden zu klein gefunden werden.

Irrt man also bei $\alpha$ kleiner Bär um $20''$, und bei jedem der beiden andern Sterne um $1''$, so kann der Fehler der Zeitbestimmung höchstens auf $1^{\mathrm{s}}1$, und der Fehler der Polhöhe höchstens auf $13''$ gehen, in so fern die Stellungen der Sterne selbst genau zuverlässig sind. Will man die möglichen Fehler der Endresultate nicht durch die Fehler der Zeiten, sondern durch die Fehler der Höhen bestimmen, so dienen dazu die andern oben gegebenen Formeln. Setzen wir nemlich, dass in den drei Beobachtungsmomenten die Höhen der Sterne um $D,D',D''$ Raum-Secunden grösser waren, als die Stellung des Sextanten erforderte, so entspringt daraus ein Fehler von
\[
+0.0430D+0.0177D'-0.0607D''
\]
Zeit-Secunden bei der Voreilung der Uhr, und von
\[
+0.446D-0{,}823D'+0.377D''
\]
Raum-Secunden bei der Polhöhe, um welche beide zu gross ausfallen werden. Die letztern Formeln dienen zugleich, den Einfluss der etwa noch bei den Stern%
\SourceLeaf{145}%
positionen Statt findenden Unsicherheiten zu schätzen; man braucht nur unter $D,D',D''$ die Fehler dieser Positionen in der Richtung der drei Vertical-Kreise und in Secunden des grössten Kreises zu verstehen, um sogleich die gefundenen Formeln für den Einfluss derselben gebrauchen zu können. Es sind nemlich $D,D',D''$ die Unterschiede der Höhen, die der Stellung des Sextanten entsprechen, von den Höhen derjenigen Punkte der Himmelskugel, welche die bei der Rechnung zum Grunde gelegten geraden Aufsteigungen und Abweichungen haben, und es ist begreiflich einerlei, ob sie bei völliger Übereinstimmung der Sterne mit diesen Punkten von den Fehlern der Beobachtungen, oder bei völlig genauen Beobachtungen von kleinen Unterschieden zwischen den Sternen und jenen Punkten herrühren. Schliesslich bemerke ich noch, dass die Summe der Coëfficienten in jeder der beiden Formeln für den Fehler der Polhöhe, so wie auch in der zweiten für den Fehler der Zeitbestimmung immer $=0$, hingegen in der ersten Formel für den Fehler der Zeitbestimmung $=1$ wird, wovon man den Grund bei einigem Nachdenken leicht finden wird.
